
\section{系统定义}
\subsubsection{1. 中文系统名称}

重庆邮电大学智慧访客预约与出入校管理系统

\subsubsection{2. 英文系统名称}

Smart Visitor Reservation and Campus Access Management System of Chongqing University of Posts and Telecommunications

\subsubsection{3. 系统设计意图}

本系统面向高校校园访客管理场景，旨在通过数据库应用系统对校外访客从预约申请、校内确认、部门审批、门岗核验、入校登记、离校登记到记录归档和统计分析的全过程进行信息化管理。系统以 MySQL 数据库为核心，对访客信息、预约信息、审批信息、通行凭证、出入校记录、黑名单、用户角色权限和系统日志等数据进行统一组织、约束、查询和统计。

\subsubsection{4. 系统目标}

建立规范化的访客预约和审批流程，减少线下沟通和纸质登记。\\
通过黑名单检查、角色权限和门岗核验提升校园出入安全性。\\
通过数据库记录保存完整审批链路和出入校轨迹，便于追溯和审计。\\
通过统计报表分析访客数量、部门接待情况、校门通行情况和超时未离校情况。\\
为数据库原理课程设计提供完整的需求分析、数据流程图、数据字典、E-R 图、关系模式、SQL 查询、系统实现和测试材料。\\

\subsubsection{5. 系统服务对象}

\begin{longtable}{p{0.47\textwidth}p{0.47\textwidth}}
\toprule
服务对象 & 说明 \\
\midrule
访客 & 校外来访人员，使用系统提交预约申请并查看预约状态和通行凭证 \\
被访人 & 校内教师或工作人员，负责确认访客预约是否属实 \\
部门审批人员 & 对本部门相关访问申请进行审核 \\
门岗安保人员 & 在校门处核验访客通行凭证并办理入校、离校登记 \\
系统管理员 & 维护用户、角色、权限、部门、校门、字典、黑名单和日志 \\
校级管理人员 & 查看全校访客记录和统计报表，为管理决策提供依据 \\
\bottomrule
\end{longtable}


\subsubsection{6. 系统核心功能}

访客预约申请：登记访客身份、联系方式、来访事由、被访人、访问时间、车辆和随行人员信息。\\
黑名单检查：在提交预约和门岗核验环节检查访客是否存在限制入校记录。\\
被访人确认：被访人对预约申请进行确认或拒绝。\\
部门审批：部门审批人员对已确认预约进行审批。\\
通行凭证生成：审批通过后生成通行码并记录有效期和使用状态。\\
门岗核验：门岗安保人员核验通行码、访客身份和预约状态。\\
入校登记与离校登记：记录访客进出校门、时间、经办人员和访问状态。\\
当前在校与超时未离校管理：查询已入校未离校访客并识别超时风险。\\
黑名单管理：维护黑名单原因、限制期限和处理状态。\\
访客记录查询与统计报表：支持按访客、部门、校门、时间、状态等条件查询统计。\\
系统管理：维护用户、角色权限、部门、校门、字典、通知和系统日志。\\
自动截图与报告生成：为课程设计报告自动生成运行截图和文档素材。\\

\subsubsection{7. 系统建设意义}

从校园管理角度看，本系统能够提升访客预约审批效率，强化出入校安全控制，减少人工登记差错。从数据库课程设计角度看，本系统业务流程完整、数据对象丰富、实体关系清晰，适合展示数据库需求分析、概念结构设计、逻辑结构设计、SQL 查询设计和系统实现测试等核心能力。

\subsection{本章小结}
本系统围绕高校访客预约和出入校管理构建统一数据库应用，系统名称、业务范围、服务对象和建设意义已经明确，为后续需求分析、概念结构设计和系统实现提供边界。
